The supplied old controller and pre-synthesized new controller are fixed. Baseline PC2's old requirements and \ContractPCTwoNew{} new requirements remain; added new-endpoint calibration availability yields \PCTwoTraceNewStarts{} new starts, separate from the interval duty. Transfers map old destinations to calibration counterparts, retaining ordinary observer memory. Work, calibration requests and updates are controllable; arrivals, tool outcomes, resets and calibration completion are UC. Entry inherits old memory and initializes the interval monitor; new-start looks up the current ordinary observer tuple. No precedence is supplied. Targets match components and new-monitor memory at reachable new-endpoint states.

Interval $R_{CAL\_READY}$ observes transfers as calibration entry; endpoint availability observes ordinary requests. Completion restores both. Handover matches component, observer and new-monitor memory; real-history and runtime obligations remain unvalidated.


\paragraph{Generation responsibilities beyond Cell.}\label{ta:pc2-generation}
PC2 separates an application obligation from its generated update protocol. The designer supplies the two old/new arm models, their local transfer maps, twelve old and thirteen new requirements, the interval duty and fixed endpoints. For interval availability, each readiness bit starts true, becomes false on its arm's transfer or calibration request, and becomes true on completion; their disjunction must hold. This obligation is supplied, not inferred from the plant.

From these inputs, FG-DUCS constructs the twenty-five independent requirement boundaries and two local transfers, their active-monitor updates, mixed-version successors, UC-quiescence guards and target tests. The frontend prepares current-observer lookup maps and physical closure; the lazy solver generates the update product on demand. The saved result covers 126 entries with 381 policy states, 522 edges and maximum rank 37. Thus neither the sequence nor those policy states are part of the input. These counts describe one saved result, not a bound on the full mixed product.

For the published DUCS interface, a representation of these independently changing arms and requirement memories would place their staged behavior and admissibility in the supplied $E'$ and transition requirements. DUCS itself still generates its global update protocol; that work is not assigned to the designer~\cite[Secs.~4.2, 5]{Nahabedian2020}. A GR(1) representation would supply corresponding version, local-state, active-memory and observer variables with transition constraints; its existing construction generates switching bookkeeping and the bridge~\cite[Def.~3.1, Sec.~4]{Amram2022GR1DynamicUpdate}. Unlike the deterministic Cell construction, PC2 has UC outcomes, so this responsibility comparison does not establish an equivalent GR(1) game or a general translation. It identifies which local operations FG-DUCS generates directly, without measuring modeling effort, minimum encoding size or external-solver performance.
